\section{讨论、有效性威胁与伦理影响}
\label{sec:discussion}

\subsection{测量效度}

两个默认目标覆盖的行为范围都很窄。KeywordRate 只检测显式子串，无法识别委婉拒答、答非所问或无标记的安全回避，也可能误伤对拒答话术的正常讨论。默认英文标记对中文回答尤其不足；主案例虽加入中文标记，仍没有提供语义分类器或人工标注。首词元 KL 则只比较一个位置的词表分布，无法代表后续序列、事实正确性、长程推理或工具调用。两个分值适合作为快速搜索代理，但不足以支撑“总体能力无损”或“模型质量提高”。

单方向假设也具有范围边界。式~\eqref{eq:direction} 用两组均值差压缩复杂行为，便于形成低秩编辑；Piras 等人却表明，多方向方法可在其设置中比单方向基线更有效~\cite{piras2026som}。Heretic 当前搜索的是方向层与作用强度，不会自动学习多个独立拒答子方向。未来实验应比较单方向、多方向与语义条件化干预，并报告不同语言、主题和提示形式上的方差。

\subsection{案例内效度与精度差异}

主案例缺少完整 journal、数据集 commit 和 BF16 成品复评分，因而无法从公开材料重放 trial 198 的全部轨迹~\cite{qwen38maincase}。4-bit 搜索会改变权重、激活与最优 Pareto 点；把 adapter 合并到 BF16 后，即使参数相同，代理分值也可能漂移。严谨复现应在最终 BF16 工件上重新计算关键词率与首词元 KL，并增加语义拒答、完整序列分布、标准任务和安全评估。只有公开原始逐 trial 记录时，才适合绘制 Pareto 散点或统计搜索稳定性。

thinking 模板是另一项内效度威胁。开放 \code{<think>} 位于提示端还是新生成端，决定共同前缀探测能否观察到它。错误切点会让方向和 KL 主要描述思考开场，而不是 instruct 回答。每次运行都应保存最终渲染提示、精确 \code{response\_prefix}、模板开关和测量位置，并分别评估 thinking 与 non-thinking 行为。

\subsection{架构与版本外推}

源码映射能解释为何相同属性形状可覆盖 Gated Attention 与 Gated DeltaNet，却不能证明 Qwen3.8 的所有辅助模块都已处理。文本 Prompt 不经过图像输入，视觉编码器与多 token 预测权重没有消融，当前测试矩阵也没有 Qwen3.8-27B。纯状态空间架构、属性名变化或新的投影包装都可能使模块发现失败。后续工作需要固定小型 Qwen3.8 架构夹具，对 64 层组件计数、LoRA target、前缀渲染和导出加载做集成测试，再单独评价视觉与视频路径。

版本冻结仍存在两个代码边界。第一，模型类别探测读取远程 config 时没有传入 \code{model\_commit}，随后 tokenizer 与模型本体才使用该 revision。第二，创建复现目录时，代码会查询当前 Hub 模型信息并覆写 \code{settings.model\_commit}；若实际运行显式固定了旧 revision，生成材料可能记录成生成当时的 Hub HEAD。据此推断，复现包应同时保存“实际加载 revision”和“生成时查询 revision”，并在不一致时中止或明确报警。

\subsection{双重用途与责任边界}

方向消融可能减少合法主题上的过度拒答，也可能削弱针对危险请求的恰当拒答。公开机制分析有助于研究对齐脆弱性、红队评估和可解释性，但同一技术也能扩大危险能力暴露。使用者应遵守模型与数据许可、机构伦理规则和适用法律，在隔离环境中评估修改后的模型，并保留修改来源、参数与分值记录。本文不提供绕过在线服务控制、隐藏权重来源或部署滥用系统的方法；任何发布决定都应基于下游安全评估，而不是单一拒答分数。
